\section{Experimenty}

\subsection*{Hledání hyperparametrů}

Hyperparametry jsou hledány metodou random search~\cite{bergstra2012random}. Oproti mřížkovému
hledání vyzkouší při stejném počtu běhů více různých hodnot každého hyperparametru, což je
výhodné, pokud výsledek ovlivňuje jen několik z nich. Hledání probíhá zvlášť pro model A a model B,
protože se optimální hodnoty (zejména learning rate) mezi encoderem trénovaným od nuly a malou
hlavou nad zmrazeným encoderem mohou výrazně lišit.

Prohledávaný prostor:
\begin{itemize}
    \item learning rate $\eta$: log-uniformně z intervalu $[10^{-4}, 3 \cdot 10^{-3}]$,
    \item weight decay $\lambda$: log-uniformně z intervalu $[10^{-5}, 10^{-2}]$,
    \item velikost dávky: $\{32, 64, 128\}$,
    \item šířka skryté vrstvy hlavy: $\{64, 128, 256\}$.
\end{itemize}
Learning rate a weight decay se vzorkují na logaritmické škále, protože jejich vliv závisí
na řádu hodnoty, ne na absolutním rozdílu. Každá konfigurace je natrénována s jedním seedem
a vyhodnocena podle MAE na validační množině. Pro každý model je vyzkoušeno 8 konfigurací.
Konfigurace se vzorkují se stejným seedem, takže oba modely dostávají stejné kombinace
hyperparametrů. Nejlepší konfigurace každého modelu je poté natrénována znovu se třemi seedy.

\begin{table}[h]
    \centering
    \begin{tabular}{crrrrrr}
        \toprule
        & $\eta$ & $\lambda$ & Dávka & $h$ & MAE A & MAE B \\
        \midrule
        1 & $1{,}02 \cdot 10^{-3}$ & $1{,}45 \cdot 10^{-5}$ & 128 & 64 & \textbf{1,1561} & 1,1407 \\
        2 & $1{,}87 \cdot 10^{-4}$ & $3{,}37 \cdot 10^{-5}$ & 64 & 256 & 1,1714 & 1,1425 \\
        3 & $2{,}31 \cdot 10^{-3}$ & $6{,}76 \cdot 10^{-5}$ & 128 & 256 & 1,1805 & 1,1538 \\
        4 & $2{,}06 \cdot 10^{-3}$ & $3{,}46 \cdot 10^{-4}$ & 32 & 64 & 1,1607 & 1,1450 \\
        5 & $1{,}65 \cdot 10^{-3}$ & $4{,}38 \cdot 10^{-5}$ & 64 & 256 & 1,1742 & 1,1474 \\
        6 & $8{,}52 \cdot 10^{-4}$ & $6{,}06 \cdot 10^{-3}$ & 128 & 64 & 1,1587 & \textbf{1,1393} \\
        7 & $1{,}51 \cdot 10^{-3}$ & $3{,}59 \cdot 10^{-4}$ & 128 & 64 & 1,1614 & 1,1443 \\
        8 & $1{,}76 \cdot 10^{-4}$ & $3{,}11 \cdot 10^{-4}$ & 32 & 128 & 1,1606 & 1,1521 \\
        \bottomrule
    \end{tabular}
    \caption{Random search: validační MAE osmi konfigurací (seed 123). Tučně nejlepší
    konfigurace každého modelu.}
    \label{tab:search}
\end{table}

Výsledky searche jsou v tabulce~\ref{tab:search}. Nejlepší je pro model A konfigurace~1 a pro
model B konfigurace~6, obě s velikostí dávky 128 a šířkou hlavy $h = 64$. Rozdíly mezi
konfiguracemi jsou malé: validační MAE se pohybuje mezi 1,156 a 1,181 u modelu A a mezi 1,139
a 1,154 u modelu B. U modelu A dosáhly nejvyšší chyby všechny tři konfigurace se šířkou hlavy 256.

\subsection*{Experiment 1: přínos jednotlivých modalit}

\textbf{Hypotéza:} Obrázky nesou informaci o popularitě, která v metadatech není. Model
kombinující obrázky s metadaty proto dosáhne nižší chyby než model nad samotnými metadaty.
Samotné obrázky budou méně informativní než metadata.

\textbf{Postup:} Porovnávají se tři varianty modelu B: pouze metadata, pouze obrázky a obě
modality. Varianta pouze s metadaty je samotná hlava $g_\phi$ bez encoderu, varianta pouze
s obrázky používá encoder ConvNeXt-Tiny a do hlavy vstupuje jen vektor obrázků. Všechny varianty
sdílejí stejné rozdělení dat, trénovací postup a hyperparametry nejlepší konfigurace modelu B. Každá varianta je natrénována se třemi seedy a vyhodnocena
pomocí MAE a RMSE na validační množině.

\subsection*{Experiment 2: encoder CLIP}

\textbf{Hypotéza:} Model CLIP~\cite{radford2021clip} je předtrénován kontrastivně na dvojicích
obrázků a textů z internetu, ne na klasifikaci do tříd ImageNetu. Jeho příznaky proto lépe
popisují vlastnosti propagačních obrázků her, jako je styl, žánr nebo kvalita zpracování.
Model B s encoderem CLIP dosáhne nižší chyby než s encoderem ConvNeXt-Tiny.

\textbf{Postup:} V modelu B je encoder ConvNeXt-Tiny nahrazen obrazovou částí modelu CLIP
ViT-B/32 s vahami od OpenAI z knihovny OpenCLIP~\cite{ilharco2021openclip}. Encoder je také
zmrazený a jeho výstupem je vektor délky $d = 512$. Obrázky $128 \times 128$ se zvětší
na $224 \times 224$, na kterých byl CLIP předtrénován, a normalizují se jeho průměrem
a směrodatnou odchylkou. Hlava, trénovací postup a hyperparametry nejlepší konfigurace modelu B
se nemění. Hyperparametry byly vybrány pro encoder ConvNeXt, což může CLIP mírně znevýhodnit.
Model je natrénován se třemi seedy a vyhodnocen pomocí MAE a RMSE na validační množině.
